\section{Заключение}

В ходе выполнения лабораторной работы были исследованы вольтамперные характеристики полупроводникового диода и работа однополупериодного выпрямителя. По полученным результатам можно сделать следующие выводы:
\begin{enumerate}
    \item На рисунке \ref{fig:diode_vah} можно наблюдать, что у диода есть пороговое напряжение, начиная с которого он проводит ток. Это достигается в момент, когда внешнее электрическое поле преодолевает потенциальный барьер на p-n переходе диода, что позволяет создавать ток.
    \item На том же рисунке \ref{fig:diode_vah} ещё можно заметить, что динамическое сопротивление диода всегда меньше статического сопротивления. Это можно объяснить математически: динамическое сопротивление диода определяется как производная тока по напряжению, а статическое сопротивление --- как отношение напряжения к току. Так как ток в диоде растёт экспоненциально с ростом напряжения, то производная тока по напряжению всегда будет больше, чем отношение напряжения к току. Другими словами, тангенс угла наклона касательной к кривой ВАХ диода всегда будет больше, чем тангенс угла наклона секущей, проведённой через начало координат и точку на кривой ВАХ диода.
    \item Максимальное значение напряжения на нагрузке меньше амплитуды источника питания. Это показывает, что диод не является идеальным ключом, и часть входного напряжения падает на p-n переходе.
    \item Максимальное обратное напряжение на диоде практически равно амплитуде источника питания. Это означает, что в схеме однополупериодного выпрямителя к закрытому диоду прикладывается почти всё напряжение источника. Это важно учитывать при выборе диода для выпрямителя, так как он должен быть способен выдержать обратное напряжение.  
\end{enumerate}